\textbf{Zusammenfassung}: Was ist Metadatenaktivismus? Dieser Beitrag
skizziert Erschließungsarbeit, Motivation, Prinzipien und Perspektiven
für die Erschließung von virtuellen Ausstellungen, die seit 2019 mit
`DDBstudio', dem Ausstellungsservice der Deutschen Digitalen Bibliothek
(DDB) entstehen. Mittelkürzungen stellen diesen Dienst der DDB in Frage.
Ausgehend vom Netzwerk `WikiKult -- Offene Kulturdaten', einem freien
Zusammenschluss von Expert:innen aus dem Kulturerbesektor, werden
Ausstellungen der DDB seit Februar 2026 in Wikidata und in
Bibliothekssystemen erschlossen. Angesichts der wachsenden Kooperation
formulieren die Autoren nicht zuletzt das Ziel, dass Bibliotheken,
Archive und andere Kulturinstitutionen, insbesondere die Deutsche
Digitale Bibliothek und die Wikimedia-Bewegung, anhand und für
wirkungsvolle Projekte öfter zusammenarbeiten -- zum Beispiel für neue
virtuelle Ausstellungen über 2026 hinaus.

\textbf{Abstract}: What is metadata activism? This article outlines the
work of cataloging, the motivation, principles, as well as perspectives
for cataloging virtual exhibitions created since 2019 with DDBstudio, an
exhibition service of the German Digital Library (DDB). Budget cuts are
jeopardizing this DDB service. Based on the `WikiKult -- Open Cultural
Data' network, an independent association of experts from the cultural
heritage sector, DDB exhibitions have been cataloged in Wikidata and
union catalogues since February 2026. In light of this growing
cooperation, the authors express the goal that libraries, archives,
other cultural institutions, like the DDB and the Wikimedia movement,
collaborate more frequently.
